\section{Introdução}

A inteligência artificial generativa tem ampliado a demanda por infraestrutura computacional capaz de processar grandes volumes de dados e executar modelos com elevado custo de treinamento ou inferência. Essa demanda aproxima a área de IA dos desafios tradicionalmente estudados em computação de alto desempenho, como uso eficiente de recursos, gargalos de memória, paralelismo, aceleração por hardware e escalabilidade.

Mesmo quando modelos de menor escala são utilizados em ambientes educacionais, sua execução evidencia limitações associadas a processamento, memória, armazenamento, comunicação entre componentes e, quando disponível, aceleração por GPU. Assim, experiências práticas em laboratório permitem observar, ainda que em escala reduzida, aspectos que também aparecem em ambientes profissionais de IA e HPC.

Este trabalho investiga como características de infraestrutura de hardware influenciam a execução de cargas relacionadas à IA generativa em um ambiente de laboratório. A pergunta de pesquisa que orienta o estudo é: \textit{como características de infraestrutura de hardware observadas em laboratório influenciam o desempenho e a viabilidade de execução de cargas de IA generativa em um contexto introdutório de HPC?}

O objetivo geral é analisar, com base em experimentos e medições técnicas, a relação entre componentes de hardware e a execução de cargas associadas à IA generativa. Como objetivos específicos, busca-se caracterizar o ambiente de execução, coletar métricas de desempenho e utilização de recursos, identificar gargalos e discutir os resultados à luz de conceitos de arquitetura de computadores e HPC.

A principal contribuição esperada é uma análise metodológica e experimental adequada a um contexto de iniciação científica, conectando os conhecimentos construídos em laboratório com um problema atual da área de sistemas computacionais de alto desempenho.

O restante do artigo está organizado da seguinte forma. A Seção~\ref{sec:fundamentacao} apresenta os conceitos fundamentais. A Seção~\ref{sec:relacionados} discute trabalhos relacionados. A Seção~\ref{sec:metodologia} descreve a metodologia. A Seção~\ref{sec:resultados} apresenta os resultados obtidos. A Seção~\ref{sec:discussao} discute os achados e limitações. Por fim, a Seção~\ref{sec:conclusao} apresenta as conclusões e trabalhos futuros.
